\section{Introduction}
\label{sec:intro}

Commercial buildings and workplace electric-vehicle (EV) fleets are converging
on a shared electrical envelope. When chargers are bidirectional, parked
vehicles become schedulable storage for the building---vehicle-to-building
(V2B) coordination---and the design questions multiply: how many chargers does
a site need, how much peak reduction can a fleet deliver, how should charging
be scheduled against a time-of-use tariff, and how do these answers change
with climate, building type, or workforce composition. Methods that answer
such questions, whether optimization-based or learned, must be developed and
evaluated on data that couples several channels at once: the building load,
the charging sessions, the drivers' arrival and departure behavior, the
tariff, demand-response (DR) events, and any on-site photovoltaics (PV) or
stationary storage, all on a common clock.

That data largely does not exist. Three barriers recur. First,
\emph{privacy}: charging records identify drivers and facilities, so
operational logs rarely leave the utility or campus that collected them.
Second, \emph{site lock-in}: the public exceptions---ACN-Data
\citep{lee2019acn}, ElaadNL open data \citep{elaadnl2020}, the EV~WATTS
evaluation dataset \citep{evwatts2023}---cover a handful of venues and fixed
infrastructures; one cannot ask a Pasadena workplace dataset how a Miami
retail site would behave, nor scale it to five hundred vehicles. Third, and
least appreciated, \emph{the key control variable is unmeasured}: charging
equipment meters energy, not batteries, so no public dataset records battery
state of charge (SoC)---the quantity nearly every V2B scheduling formulation
constrains and optimizes.

Synthetic generation addresses all three barriers if---and only if---the
generated data can be trusted. A generator provides privacy by construction,
supports counterfactual scenarios and arbitrary scale, and can label
quantities (such as SoC) that measurement cannot reach. The obligation this
creates is evidential: for a synthetic dataset to support methodological
conclusions, its \emph{representativeness} must be quantified against real
data and its \emph{utility} for downstream learning must be demonstrated
rather than asserted. These are precisely the criteria the KDD Datasets and
Benchmarks track sets for data generators, and they organize this paper.

We present \generator{}, a configurable generator of coupled V2B datasets. A
scenario is specified by five descriptors---site, building, driver population,
charging equipment, and measurement noise---and one integer seed; the
generator emits a coherent bundle of fifteen-minute time series and entity
tables: building load split into flexible and inflexible components, charging
sessions, driver and vehicle attributes, chargers, tariff prices, DR events,
PV generation, and storage specifications, together with a provenance
manifest. Three properties define the design.

\begin{itemize}
\item \textbf{Physical.} Building load is not a shape template: every dataset
is produced by an EnergyPlus \citep{crawley2001energyplus} simulation of an
ASHRAE 90.1-2019 DOE prototype building \citep{pnnl_prototypes} under
typical-meteorological-year weather subjected to a controlled, per-sample
perturbation. The identical perturbed weather frame drives the PVWatts-style
PV model \citep{dobos2014pvwatts} and is exported with the dataset, so load,
PV, and weather can never disagree.
\item \textbf{Calibrated.} Driver behavior is organized on three interpretable
axes---plug-in frequency $\phifreq$, timing consistency $\kapcons$, and daily
distance $\deldist$---partitioned into behavioral regions. Arrival, dwell, and
energy-requirement distributions are fitted per region by maximum likelihood
from three real charging corpora (ACN-Data, EV~WATTS, ElaadNL), with fit
quality (Kolmogorov--Smirnov distance) and provenance stamped on every fitted
value. Because no source measures SoC, the SoC channels are \emph{explicitly
modeled priors} reconstructed from delivered energy and inferred battery
capacity; we state this contract prominently rather than presenting modeled
values as measurements.
\item \textbf{Reproducible.} Randomness is drawn from hash-keyed substreams
(one per graph node, vehicle, and day), so a scenario configuration and one
seed regenerate any dataset bitwise, on any machine, regardless of the order
in which features were added to the generator. Every run records a manifest
containing each configuration value with its source, per-file checksums, and
the validation verdict.
\end{itemize}

\paragraph{Contributions.}
(i)~A forward-sampling generative model for coupled V2B data whose behavioral
layer is calibrated per region from three real charging corpora and whose
building layer is a real EnergyPlus simulation, with a single weather
realization shared by load, PV, and the exported weather channel
(\S\ref{sec:generator}--\S\ref{sec:calibration}).
(ii)~A reproducibility contract---order-independent hash-keyed randomness,
bitwise regeneration, per-value provenance manifests---accompanied by a
datasheet \citep{gebru2021datasheets}, Croissant metadata
\citep{akhtar2024croissant}, and dual licensing (\S\ref{sec:generator},
\S\ref{sec:release}).
(iii)~A two-sided evaluation meeting the generators bar: representativeness
via ASHRAE Guideline~14 building-load fidelity against ComStock
\citep{parker2023comstock} and Building Data Genome~2 \citep{miller2020bdg2}
meters, and held-out, confidence-interval-quantified distributional fidelity
against the calibration corpora (\S\ref{sec:representativeness}); utility via
train-on-synthetic/test-on-real forecasting at parity with training on real
data on ACN-Data, plus a cross-cohort transfer study
(\S\ref{sec:utility}).
(iv)~A released benchmark suite: sixty scenario configurations spanning nine
sites, six building types, thirteen driver cohorts, and six equipment mixes;
distribution-shift cohort pairs; and an 18{,}000-sample ten-building reference
corpus (\S\ref{sec:release}).
